\subsubsection{Step Budget}\label{sec:res:fewstep}

The flow-based models keep their success with constraint satisfaction as the step budget falls to a
single network evaluation, while the diffusion baseline does not reach it at a small budget at all
(\autoref{fig:k-ladder}). The baseline reaches $0.667$ at $\nfe=1$ and $1.000$ at $\nfe=10$ and
$\nfe=20$. Over the same range the flow-based models never leave the top of the axis: both
average-velocity models read $1.000$ or $0.967$ at one and two evaluations, analytic average-velocity
matching holds $0.933$ out to ten, and instantaneous-velocity matching reads $1.000$ at one, two and
twenty.

The two curves are not read the same way, and the difference is the point. Each diffusion point is its own
model: the number of denoising steps is fixed when the noise schedule is discretised during training, so
the budgets of that curve are three separately trained checkpoints, one per step count, and moving along
it costs a training run. The flow-based curve is one checkpoint per model, read at every budget, because
there the step count is the number of steps the ODE solver takes at inference. For a flow-based model the
budget is therefore a dial that can be turned after training, and for the diffusion baseline it is a
property of the checkpoint --- the same asymmetry \autoref{fig:raw-plans} shows, where the two diffusion
panels are two separately trained models and each flow-based row is one model read twice. The cost of
that asymmetry is a training run: the $\nfe=2$ panel of that figure required one, and the single-seed
evaluation it produced puts the baseline at about $225$\,ms per control step at two denoising steps,
against instantaneous-velocity matching's $17.3$\,ms at one. Cutting the budget from twenty to two
leaves the baseline at thirteen times the cost.

\begin{figure}[htpb]
  \centering
  \includegraphics[width=\textwidth]{fig_avoiding_k_ladder}
  \caption[Success against step budget]{Success with constraint satisfaction against the step budget,
  averaged over the three geometries; the protocol of \ac{DPCC}, five training seeds $\times$ two
  episodes per geometry, temporal U-Net of $4.0$\,M parameters. Each model is drawn at the selection
  rule it is operated with: temporal consistency for the two average-velocity models, cumulative
  projection cost for instantaneous-velocity matching and for the baseline. Each is drawn at the budgets
  it was evaluated at, and every point carries all five seeds and all three geometries.
  Consistency-interpolated average-velocity matching lies on instantaneous-velocity matching at both of
  its budgets, so it is dashed and drawn with the smaller marker. Legend: MeanFM, CI-MeanFM, FM and the
  diffusion model of \ac{DPCC}. CI-MeanFM and FM are both at $1.000$ at both of their budgets, so their lines coincide; they are drawn
  displaced by $\pm0.006$ on the outcome axis so that both are visible. The displacement is a drawing
  offset and not a difference in the data.}
  \label{fig:k-ladder}
\end{figure}

All four series are the same sample --- five training seeds, two episodes, three geometries --- so the
figure is one comparison and not several. What separates the baseline's curve from the other three is
what a point on it costs to produce: its three budgets are three separately trained checkpoints, and the
flow-based curves are one checkpoint per model read at each budget. $\nfe$ is an inference-time
accuracy--cost dial for the flow-based ODEs and a training decision for the diffusion baseline, which a
later choice cannot undo.
